\section{Алгоритмы распознавания языка}
\label{sec:algorithms}

По варианту~13 нужно понять, на каком языке написан PDF: на русском или на немецком. Для этого сравниваются три способа \mdash короткие слова, частотные слова и нейросеть.

Сначала по учебным текстам для каждого языка собирается поисковый образ языка (ПОЯ): какие слова или куски текста для него обычны. У нового документа по тем же правилам строится поисковый образ документа (ПОД). Дальше смотрим, к какому языку документ ближе. Для слов это произведение вероятностей, для нейросети \mdash вероятность класса. Дополнительно считается расстояние Out-of-Place: насколько по-разному выстроены списки частых элементов.

Учебные тексты лежат в каталоге \texttt{corpus/train}: по семь медицинских заметок на каждый язык, примерно 50~Кбайт русского текста и 85~Кбайт немецкого. Для проверки отдельно приготовлены десять коротких фраз с уже известным языком. В обучение они не входят.

Перед сравнением текст приводится к нижнему регистру, буква <<ё>> заменяется на <<е>>. Немецкие умлауты сохраняются. Русские слова приводятся к начальной форме, немецкие только к нижнему регистру. Короткие служебные слова вроде <<и>> или <<der>> не выбрасываются: по ним языки хорошо отличаются.

В методе коротких слов в образ языка попадают слова не длиннее пяти букв, если в учебных текстах они встретились больше трёх раз. Чем слово чаще, тем выше его вероятность. Если слова в образе нет, ему даётся очень маленькая вероятность, чтобы расчёт не обрывался. Документ оценивается как произведение вероятностей его коротких слов:

\begin{equation}
	P(\mathrm{lang}\mid\mathrm{doc}) = \prod_{l \in L_{\mathrm{short}}} P(l\mid\mathrm{lang}).
	\label{eq:short}
\end{equation}

Берётся язык с большим значением. Считают сумму логарифмов: само произведение на длинном тексте слишком быстро становится нулём.

В методе частотных слов в образ входят пятьдесят самых обычных слов, длина здесь не важна. С документом сравниваются только слова из этого короткого списка. Формула та же, что и~\eqref{eq:short}, поэтому способ обычно быстрее.

Нейросеть смотрит не на целые слова, а на короткие куски текста длиной от одного до пяти символов: такой кусок в тексте есть или нет. Она учится на тех же заметках и в конце выдаёт две вероятности \mdash русский и немецкий. Расстояние Out-of-Place для неё не считается.

Out-of-Place сравнивает два списка, упорядоченных от частых элементов к редким. Для каждого элемента смотрят, насколько разъехались его места в списках. Если элемента во втором списке нет, разница берётся большой:

\begin{equation}
	d(A,B) = \sum_{g \in A \cup B} \lvert \mathrm{pos}_A(g) - \mathrm{pos}_B(g) \rvert.
	\label{eq:oop}
\end{equation}

Число растёт вместе с длиной списков, поэтому у коротких слов и у пятидесяти частых слов оно разного масштаба. Язык словесные методы всё равно выбирают по вероятности~\eqref{eq:short}.
